% Generated draft: requires mathematical and source review.
\documentclass[11pt]{article}
\usepackage[margin=25mm]{geometry}
\usepackage{fontspec}
\setmainfont{Latin Modern Roman}
\usepackage{amsmath,amssymb,mathtools,unicode-math}
\setmathfont{Latin Modern Math}
\usepackage{enumitem,needspace,xurl,hyperref}
\setlength{\parindent}{0pt}
\setlength{\parskip}{5pt}
\setlength{\emergencystretch}{3em}
\newcommand{\N}{\mathbb{N}}
\newcommand{\Z}{\mathbb{Z}}
\newcommand{\Q}{\mathbb{Q}}
\newcommand{\R}{\mathbb{R}}
\newcommand{\C}{\mathbb{C}}
\newcommand{\F}{\mathbb{F}}
\newcommand{\A}{\mathbb{A}}
\newcommand{\PP}{\mathbb{P}}
\newcommand{\E}{\mathbb{E}}
\newcommand{\eps}{\varepsilon}
\newcommand{\dd}{\,\mathrm{d}}
\begin{document}
\section*{UnsolvedMath ID: 30002300}
\textbf{Source problem code:} catalogue entry
\subsection{Definitions and mathematical statement}
For a positive integer \(d\), let \(P\subseteq\R^d\) be a compact polyhedron with vertex set \(V\). Write \(\operatorname{aff}(S)\) for the affine hull of a set \(S\). The hypothesis under discussion is
\[
H_d(V):\qquad \operatorname{aff}(S)=\R^d\quad\text{for every }S\subseteq V\text{ with }|S|=d+2.
\]
This condition is vacuous when \(|V|<d+2\). Define the collection of nondegenerate simplices supported on vertices of \(P\) by
\[
\mathcal S_d(V)=\bigl\{\operatorname{conv}(T):T\subseteq V,\ |T|=d+1,\ \operatorname{aff}(T)=\R^d\bigr\}.
\]
Here \(\operatorname{conv}(T)\) denotes the convex hull of \(T\).

To make the measure identity explicit in the full-dimensional interpretation, let \(\lambda_d\) be Lebesgue measure on \(\R^d\), suppose \(\lambda_d(P)>0\), and define the uniform probability measure on any measurable set \(Q\) of finite positive volume by
\[
\mu_Q(A)=\frac{\lambda_d(A\cap Q)}{\lambda_d(Q)}\qquad(A\subseteq\R^d\text{ Borel}).
\]
A signed simplex representation means that there exist an integer \(m\geq1\), simplices \(\Delta_1,\ldots,\Delta_m\in\mathcal S_d(V)\), and coefficients \(c_1,\ldots,c_m\in\R\) such that
\[
\mu_P=\sum_{j=1}^{m}c_j\mu_{\Delta_j}
\]
as finite signed Borel measures. The coefficients may be negative; this normalization forces \(\sum_{j=1}^{m}c_j=1\). If uniform measure instead means the unnormalized restriction \(\lambda_d\!\restriction_Q\), the corresponding representation uses that restriction on each set. For sets of positive \(d\)-dimensional volume, the existence of a representation is equivalent under these two conventions by rescaling coefficients.

\textbf{Problem (supplied source formulation).} Can the hypothesis \(H_d(V)\) in the signed simplex representation theorem be omitted while retaining the stated representation for every polyhedron in the theorem's intended class? If it cannot be omitted, can it be replaced by a strictly weaker sufficient hypothesis on \(V\)? No particular replacement hypothesis is specified.

\textbf{Unresolved scope qualifications.} The supplied record does not reproduce the complete theorem. In particular, it does not specify whether ``compact polyhedron'' means a convex polytope or a broader polyhedral set, whether full dimensionality is required, or whether simplices must satisfy additional restrictions such as containment in \(P\). The measure identity above is an explicit full-dimensional interpretation of the supplied conclusion, not a reconstruction of missing hypotheses. If lower-dimensional polyhedra or degenerate simplices are intended, the reference measure and admissible simplex dimensions must be specified. The original theorem must therefore be consulted before this interpretation can serve as a definitive statement.
\subsection{Short English statement}
The question asks whether the condition that every \(d+2\) vertices affinely span \(\R^d\) can be removed or weakened while still expressing a compact polyhedron's uniform measure as a finite real linear combination of uniform measures on simplices with \(d+1\) vertices from its vertex set. The exact class of polyhedra and the complete theorem hypotheses remain unspecified in the supplied record.
\subsection{Sources}
\begin{itemize}
\item Local UnsolvedMath registry record: permanent ID \(30002300\), source problem code \texttt{catalogue entry}, in the supplied project registry context \texttt{data/unsolvedmath/problems.json} and supplementary \texttt{catalog/problems.json}. Title: \emph{Weakening General Position in Signed Simplex Representations}. The supplied collection is \emph{Oberwolfach Reports Open Problems}; the proposed year is 2013, and no individual proposer is identified.
\item Supplied primary reference: \emph{Structured Function Systems and Applications} (2013), Oberwolfach Reports, pp.\ 579--655, DOI 10.4171/owr/2013/11. \url{https://doi.org/10.4171/owr/2013/11}. The record describes Oberwolfach Reports as documenting mathematical workshops and their open-problem discussions.
\item Supplied literature source: ``Oberwolfach Report 11/2013.'' \url{https://doi.org/10.4171/owr/2013/11}.
\item Supplied literature source: ``Scheiderer 2006.'' \url{https://doi.org/10.1007/s00229-006-0630-5}. No fuller bibliographic description or specific supporting result is supplied.
\item Supplied literature assessment, checked on \texttt{2026-08-21}: the record labels the problem open and says that the report explicitly asks about removing or weakening the affine-spanning hypothesis. It further says that subsequent checked sources give related positive and negative sum-of-squares results without settling this formulation. This is supplied dated curation derived from the OWR record, not an independently verified assessment; the record itself requests independent verification of time-sensitive status and cited claims.
\item Supplied statement-verification metadata labels the statement \texttt{corrected\_verified} and says that the official report supplies the signed simplex representation omitted from the extracted theorem reference. This verification claim is reproduced as source metadata and has not been independently checked here.
\end{itemize}
\end{document}
